\section{引言}

风电已成为现代电力系统中的重要可再生能源 \cite{ref1,ref2}。然而，风电出力受天气过程、地形条件和场站运行状态共同影响，具有显著的间歇性和波动性，会增加调度计划、备用容量配置和实时平衡的难度。因此，准确的风电功率预测是降低风电不确定性影响、支撑电力系统安全经济运行的重要工具。

随着风电场以区域场群形式接入电网，预测任务不再局限于单个场站出力，而可以面向目标区域内多个风电场站的聚合出力。本文关注区域集群总功率预测，即在给定目标窗口前可用观测和未来气象预报的条件下，预测研究区域内风电场站的总出力序列。该任务中的近期功率观测不仅反映单个场站的短期惯性，还包含多场站共同运行状态和场站间耦合关系；未来网格化气象预报则提供覆盖研究区域的外部气象驱动。

已有研究已经从多个方向提升了场群级风电预测性能，包括 CNN-LSTM 混合结构、网格化表示、图神经网络和趋势感知多步模型等 \cite{ref1,ref2,ref3,ref4,ref5,ref6,ref7,ref8,ref9,ref10,ref11,ref12}。同时，网格化数值天气预报和深度气象预报模型为未来窗口提供了重要外生信息，已有融合框架通常通过特征拼接、并行编码器或注意力交互将气象信息与功率信息结合 \cite{ref13,ref14,ref15,ref16,ref17,ref18}，也有工作直接对网格化气象场进行编码 \cite{ref19,ref20,ref21}。本文使用盘古气象大模型（Pangu-Weather）\cite{ref22} 作为预报来源。尽管这些方法证明了气象预报和功率观测融合的有效性，但多数框架仍倾向于将二者作为相对对称的输入分支处理，没有充分体现二者在业务化日内预测中的不同角色。

本文的核心出发点是：在本文的业务化日内集群总功率预测设定中，未来气象预报和近期场站功率承担不同角色。未来气象预报覆盖整个预测窗口，是驱动未来风电变化的主要外生因素；近期场站功率只在目标窗口前可用，主要表征当前运行状态、场站间耦合和短期惯性。因此，二者不应仅通过对称拼接或并行编码进行融合，而应采用更明确的角色划分：以未来气象预报作为主要驱动，以近期功率动态作为状态条件，对未来气象驱动的读取过程进行条件化约束。

基于这一认识，本文提出一种功率条件化气象记忆网络。该模型首先将历史和未来网格化气象场编码为可被预测查询读取的气象记忆；随后，将近期场站级功率动态映射为样本特定的 dynamic power context tokens，并注入气象记忆空间；最后，由并行 future queries 读取功率条件化后的气象记忆，输出 24 步集群总功率预测结果。该设计实现了跨模态的非对称条件化融合，使模型在读取未来气象驱动时显式受到当前功率状态的约束。

本文在吉林区域风电场群数据集上验证所提出的框架。实验在严格日样本和滑窗协议下进行，预测目标为区域集群总功率。结果表明，所提出模型在整体精度上优于 Weather-only 和多种网格气象编码基线；逐步长分析进一步显示，近期功率上下文对预测窗口前段具有更明显的补充作用；消融实验验证了真实功率状态、功率动态编码和气象记忆结构的必要性。

本文的主要贡献如下：

\begin{enumerate}
\item 本文从业务化区域集群总功率预测的时间可用性出发，明确区分未来气象预报和近期场站功率的建模角色：前者作为预测窗口内的主要外生驱动，后者作为目标窗口前的当前运行状态和多场站耦合信息。该视角为后续的非对称跨模态融合提供了建模依据。
\item 本文提出了一种功率条件化气象记忆网络，将历史和未来网格化气象场编码为可被预测查询读取的气象记忆，并将近期场站级功率动态转换为样本特定的 dynamic power context tokens 注入气象记忆空间，实现面向未来气象驱动的状态条件化读取。
\item 本文在吉林区域风电场群集群总功率预测任务上进行了系统验证。实验在严格时间划分和一致评价设置下比较 Weather-only、CNN、3D CNN、ViT、Swin 和 Video Transformer 等基线，并通过消融实验、误差分布和逐预测步分析验证了真实功率状态、功率动态编码和气象记忆结构的有效性。
\end{enumerate}
